This study evaluates one-hour-ahead Bitcoin volatility
forecasts using five-minute BTCUSDT USD-margined perpetual
futures data from 2022 to 2025. Hourly realized variance is
forecast using a persistence benchmark, AR(24), Fixed HAR-X,
and TVP-HAR-X within a chronological
training--validation--test design.

The results show that multi-horizon volatility and current
market activity provide useful predictive information beyond
both persistence and a conventional autoregression. On the
2025 test sample, Fixed HAR-X reduces QLIKE from \(0.5054\)
for AR(24) to \(0.4764\), an improvement of approximately
\(5.75\%\). Trading activity contributes the strongest
incremental signal among the additional HAR-X predictors and
remains positively associated with next-hour volatility
throughout the validation period.

Allowing coefficients to vary over time produces a more
qualified result. TVP-HAR-X achieves the lowest test MAE and
RMSE, at \(0.5901\) and \(0.7734\), respectively. Relative to
Fixed HAR-X, these values correspond to improvements of
approximately \(1.33\%\) in absolute error and \(1.87\%\) in
squared error. Both differences are statistically significant
under the HAC loss-differential tests.

Fixed HAR-X nevertheless records the lowest test QLIKE,
at \(0.4764\), compared with \(0.4800\) for TVP-HAR-X. The
one-sided test provides no evidence that parameter variation
improves the pre-specified primary loss. Coefficient adaptation
therefore improves average log-scale accuracy but does not
deliver a uniform advantage under variance-scale evaluation.

The evidence supports two main conclusions. First,
short-horizon Bitcoin volatility is meaningfully predictable
when recent volatility is represented across several horizons
and combined with trading activity. Second, the value of
time-varying coefficients depends on the forecast objective
and market period. Fixed HAR-X provides the strongest
performance under QLIKE, while TVP-HAR-X offers modest gains
in MAE and RMSE and additional information about changes in
predictor importance.

These findings are specific to one contract, exchange, and
forecast horizon. Broader conclusions require evaluation
across additional markets, horizons, and richer
microstructure data.